\documentclass[11pt,letterpaper]{article}
\newcommand{\labid}{1-C}
\newcommand{\labtitleshort}{SPI Bring-Up}
% Shared preamble for Module 1 lab handouts.
% Usage: each handout defines \labid and \labtitleshort, then % Shared preamble for Module 1 lab handouts.
% Usage: each handout defines \labid and \labtitleshort, then % Shared preamble for Module 1 lab handouts.
% Usage: each handout defines \labid and \labtitleshort, then \input{LabHandout_Preamble}
% BEFORE \begin{document}. Compiles with a single pdflatex pass (no glossaries tool).

\usepackage[T1]{fontenc}
\usepackage[margin=1in]{geometry}
\usepackage{titlesec}
\usepackage{enumitem}
\usepackage{booktabs}
\usepackage{tabularx}
\usepackage{graphicx}
\graphicspath{{Images/}{../Images/}}
\usepackage{xcolor}
\usepackage{hyperref}
\usepackage{fancyhdr}
\usepackage{amsmath}
\usepackage{amssymb}
\usepackage{tcolorbox}
\usepackage{caption}
\tcbuselibrary{skins,breakable}

% ── Glossary (per-document, built from shared glossary.tex) ──
% Uses \makenoidxglossaries so no external makeglossaries tool is
% needed: sorting and printing is done by LaTeX itself across two
% pdflatex passes (same cadence as the TOC). Only acronyms that
% are actually \gls{}-used in the handout appear in the printed
% glossary, and each \gls{} reference is a live hyperlink that
% jumps to its entry in the Acronyms section at the end.
\usepackage[acronym,nomain,nopostdot]{glossaries}
\makenoidxglossaries
\input{../../glossary}
% Call \printhandoutglossary just before \end{document} in each handout.
\newcommand{\printhandoutglossary}{%
	\clearpage
	\printnoidxglossary[type=\acronymtype,title={Glossary}]%
}

% ── Colors (match Module 1 lesson plan) ──────────────────────
\definecolor{stevens}{RGB}{163,31,52}
\definecolor{darkgray}{RGB}{60,60,60}
\definecolor{lightgray}{RGB}{240,240,240}
\definecolor{accentblue}{RGB}{30,80,140}

\hypersetup{
	colorlinks=true,
	linkcolor=stevens,
	urlcolor=accentblue,
	citecolor=darkgray,
}

% ── Defaults for header variables (handout overrides these) ──
\providecommand{\labid}{?}
\providecommand{\labtitleshort}{Lab Handout}

% ── Running header / footer ──────────────────────────────────
\pagestyle{fancy}
\fancyhf{}
\fancyhead[L]{\small\textcolor{darkgray}{Module 1 --- Lab Handout \labid}}
\fancyhead[R]{\small\textcolor{darkgray}{\labtitleshort}}
\fancyfoot[C]{\small\thepage}
\renewcommand{\headrulewidth}{0.4pt}

% ── Section formatting ───────────────────────────────────────
\titleformat{\section}{\Large\bfseries\color{stevens}}{\thesection}{0.6em}{}
\titleformat{\subsection}{\large\bfseries\color{darkgray}}{\thesubsection}{0.6em}{}

% ── Box styles ───────────────────────────────────────────────
\tcbset{
	infobox/.style={
		enhanced, colback=blue!3, colframe=accentblue, fonttitle=\bfseries,
		title={#1}, coltitle=white, breakable, left=6pt, right=6pt, top=4pt, bottom=4pt
	},
	warning/.style={
		enhanced, colback=red!3, colframe=stevens, fonttitle=\bfseries,
		title={#1}, coltitle=white, breakable, left=6pt, right=6pt, top=4pt, bottom=4pt
	},
	deliverable/.style={
		enhanced, colback=gray!5, colframe=darkgray, fonttitle=\bfseries,
		title={#1}, coltitle=white, breakable, left=6pt, right=6pt, top=4pt, bottom=4pt
	},
}

% ── Figure helper: always draws a placeholder box (images suppressed) ─
% Usage: \labfigure{filename}{width}{caption}{label}
\newcommand{\labfigure}[4]{%
	\begin{figure}[h!]
		\centering
		\fbox{\begin{minipage}[c][2.2in][c]{#2}
			\centering\color{darkgray}\small
				\textbf{[Figure placeholder]}\\[4pt]
				\texttt{\detokenize{#1}}\\[4pt]
				\textit{#3}
		\end{minipage}}%
		\caption{#3}
		\label{#4}
	\end{figure}%
}

% ── Title banner (call at top of document body) ──────────────
\newcommand{\labheader}[2]{%
	\begin{center}
		{\LARGE\bfseries\color{stevens} Lab Handout #1}\\[3pt]
		{\Large\bfseries #2}\\[6pt]
		{\small\color{darkgray} Capstone --- Module 1: \RF{} Hardware Foundations}
	\end{center}
	\vspace{4pt}
	{\color{darkgray}\rule{\textwidth}{0.4pt}}
	\vspace{10pt}
}

% BEFORE \begin{document}. Compiles with a single pdflatex pass (no glossaries tool).

\usepackage[T1]{fontenc}
\usepackage[margin=1in]{geometry}
\usepackage{titlesec}
\usepackage{enumitem}
\usepackage{booktabs}
\usepackage{tabularx}
\usepackage{graphicx}
\graphicspath{{Images/}{../Images/}}
\usepackage{xcolor}
\usepackage{hyperref}
\usepackage{fancyhdr}
\usepackage{amsmath}
\usepackage{amssymb}
\usepackage{tcolorbox}
\usepackage{caption}
\tcbuselibrary{skins,breakable}

% ── Glossary (per-document, built from shared glossary.tex) ──
% Uses \makenoidxglossaries so no external makeglossaries tool is
% needed: sorting and printing is done by LaTeX itself across two
% pdflatex passes (same cadence as the TOC). Only acronyms that
% are actually \gls{}-used in the handout appear in the printed
% glossary, and each \gls{} reference is a live hyperlink that
% jumps to its entry in the Acronyms section at the end.
\usepackage[acronym,nomain,nopostdot]{glossaries}
\makenoidxglossaries
% Shared acronym macro container for EE800/820 lesson plan modules.
% Usage: type \IDE, \FPGA, \UART, etc. First use is expanded; later uses show acronym only.
% Requires in each document preamble:
%   \usepackage[acronym,toc]{glossaries}
%   \makeglossaries
%   \input{../glossary}
% Print used acronyms near end of document:
%   \printglossary[type=\acronymtype,title={Acronyms Used}]

% Acronym definitions (only used entries appear in the printed glossary).
\newacronym{ask}{ASK}{Amplitude Shift Keying}
\newacronym{bga}{BGA}{Ball Grid Array}
\newacronym{bom}{BOM}{Bill of Materials}
\newacronym{bram}{BRAM}{Block Random Access Memory}
\newacronym{crc}{CRC}{Cyclic Redundancy Check}
\newacronym{css}{CSS}{Chirp Spread Spectrum}
\newacronym{dma}{DMA}{Direct Memory Access}
\newacronym{dsp}{DSP}{Digital Signal Processing}
\newacronym{eda}{EDA}{Exploratory Data Analysis}
\newacronym{fifo}{FIFO}{First In, First Out}
\newacronym{fpga}{FPGA}{Field-Programmable Gate Array}
\newacronym{fsk}{FSK}{Frequency Shift Keying}
\newacronym{fsm}{FSM}{Finite State Machine}
\newacronym{gpio}{GPIO}{General-Purpose Input/Output}
\newacronym{hal}{HAL}{Hardware Abstraction Layer}
\newacronym{i2c}{I\textsuperscript{2}C}{Inter-Integrated Circuit}
\newacronym{ide}{IDE}{Integrated Development Environment}
\newacronym{ila}{ILA}{Integrated Logic Analyzer}
\newacronym{ism}{ISM}{Industrial, Scientific, and Medical}
\newacronym{isr}{ISR}{Interrupt Service Routine}
\newacronym{lut}{LUT}{Look-Up Table}
\newacronym{mcu}{MCU}{Microcontroller Unit}
\newacronym{ml}{ML}{Machine Learning}
\newacronym{ook}{OOK}{On-Off Keying}
\newacronym{per}{PER}{Packet Error Rate}
\newacronym{rf}{RF}{Radio Frequency}
\newacronym{rssi}{RSSI}{Received Signal Strength Indicator}
\newacronym{rtl}{RTL}{Register-Transfer Level}
\newacronym{snr}{SNR}{Signal-to-Noise Ratio}
\newacronym{spi}{SPI}{Serial Peripheral Interface}
\newacronym{uart}{UART}{Universal Asynchronous Receiver/Transmitter}
\newacronym{vhdl}{VHDL}{VHSIC Hardware Description Language}

% Convenience wrappers so you can type \IDE, \RF, etc.
\providecommand{\ASK}{\gls{ask}}
\providecommand{\BGA}{\gls{bga}}
\providecommand{\BOM}{\gls{bom}}
\providecommand{\BRAM}{\gls{bram}}
\providecommand{\CRC}{\gls{crc}}
\providecommand{\CSS}{\gls{css}}
\providecommand{\DMA}{\gls{dma}}
\providecommand{\DSP}{\gls{dsp}}
\providecommand{\EDA}{\gls{eda}}
\providecommand{\FIFO}{\gls{fifo}}
\providecommand{\FPGA}{\gls{fpga}}
\providecommand{\FSK}{\gls{fsk}}
\providecommand{\FSM}{\gls{fsm}}
\providecommand{\GPIO}{\gls{gpio}}
\providecommand{\HAL}{\gls{hal}}
\providecommand{\IIC}{\gls{i2c}}
\providecommand{\IDE}{\gls{ide}}
\providecommand{\ILA}{\gls{ila}}
\providecommand{\ISM}{\gls{ism}}
\providecommand{\ISR}{\gls{isr}}
\providecommand{\LUT}{\gls{lut}}
\providecommand{\MCU}{\gls{mcu}}
\providecommand{\ML}{\gls{ml}}
\providecommand{\OOK}{\gls{ook}}
\providecommand{\PER}{\gls{per}}
\providecommand{\RF}{\gls{rf}}
\providecommand{\RSSI}{\gls{rssi}}
\providecommand{\RTL}{\gls{rtl}}
\providecommand{\SNR}{\gls{snr}}
\providecommand{\SPI}{\gls{spi}}
\providecommand{\UART}{\gls{uart}}
\providecommand{\VHDL}{\gls{vhdl}}


% Call \printhandoutglossary just before \end{document} in each handout.
\newcommand{\printhandoutglossary}{%
	\clearpage
	\printnoidxglossary[type=\acronymtype,title={Glossary}]%
}

% ── Colors (match Module 1 lesson plan) ──────────────────────
\definecolor{stevens}{RGB}{163,31,52}
\definecolor{darkgray}{RGB}{60,60,60}
\definecolor{lightgray}{RGB}{240,240,240}
\definecolor{accentblue}{RGB}{30,80,140}

\hypersetup{
	colorlinks=true,
	linkcolor=stevens,
	urlcolor=accentblue,
	citecolor=darkgray,
}

% ── Defaults for header variables (handout overrides these) ──
\providecommand{\labid}{?}
\providecommand{\labtitleshort}{Lab Handout}

% ── Running header / footer ──────────────────────────────────
\pagestyle{fancy}
\fancyhf{}
\fancyhead[L]{\small\textcolor{darkgray}{Module 1 --- Lab Handout \labid}}
\fancyhead[R]{\small\textcolor{darkgray}{\labtitleshort}}
\fancyfoot[C]{\small\thepage}
\renewcommand{\headrulewidth}{0.4pt}

% ── Section formatting ───────────────────────────────────────
\titleformat{\section}{\Large\bfseries\color{stevens}}{\thesection}{0.6em}{}
\titleformat{\subsection}{\large\bfseries\color{darkgray}}{\thesubsection}{0.6em}{}

% ── Box styles ───────────────────────────────────────────────
\tcbset{
	infobox/.style={
		enhanced, colback=blue!3, colframe=accentblue, fonttitle=\bfseries,
		title={#1}, coltitle=white, breakable, left=6pt, right=6pt, top=4pt, bottom=4pt
	},
	warning/.style={
		enhanced, colback=red!3, colframe=stevens, fonttitle=\bfseries,
		title={#1}, coltitle=white, breakable, left=6pt, right=6pt, top=4pt, bottom=4pt
	},
	deliverable/.style={
		enhanced, colback=gray!5, colframe=darkgray, fonttitle=\bfseries,
		title={#1}, coltitle=white, breakable, left=6pt, right=6pt, top=4pt, bottom=4pt
	},
}

% ── Figure helper: always draws a placeholder box (images suppressed) ─
% Usage: \labfigure{filename}{width}{caption}{label}
\newcommand{\labfigure}[4]{%
	\begin{figure}[h!]
		\centering
		\fbox{\begin{minipage}[c][2.2in][c]{#2}
			\centering\color{darkgray}\small
				\textbf{[Figure placeholder]}\\[4pt]
				\texttt{\detokenize{#1}}\\[4pt]
				\textit{#3}
		\end{minipage}}%
		\caption{#3}
		\label{#4}
	\end{figure}%
}

% ── Title banner (call at top of document body) ──────────────
\newcommand{\labheader}[2]{%
	\begin{center}
		{\LARGE\bfseries\color{stevens} Lab Handout #1}\\[3pt]
		{\Large\bfseries #2}\\[6pt]
		{\small\color{darkgray} Capstone --- Module 1: \RF{} Hardware Foundations}
	\end{center}
	\vspace{4pt}
	{\color{darkgray}\rule{\textwidth}{0.4pt}}
	\vspace{10pt}
}

% BEFORE \begin{document}. Compiles with a single pdflatex pass (no glossaries tool).

\usepackage[T1]{fontenc}
\usepackage[margin=1in]{geometry}
\usepackage{titlesec}
\usepackage{enumitem}
\usepackage{booktabs}
\usepackage{tabularx}
\usepackage{graphicx}
\graphicspath{{Images/}{../Images/}}
\usepackage{xcolor}
\usepackage{hyperref}
\usepackage{fancyhdr}
\usepackage{amsmath}
\usepackage{amssymb}
\usepackage{tcolorbox}
\usepackage{caption}
\tcbuselibrary{skins,breakable}

% ── Glossary (per-document, built from shared glossary.tex) ──
% Uses \makenoidxglossaries so no external makeglossaries tool is
% needed: sorting and printing is done by LaTeX itself across two
% pdflatex passes (same cadence as the TOC). Only acronyms that
% are actually \gls{}-used in the handout appear in the printed
% glossary, and each \gls{} reference is a live hyperlink that
% jumps to its entry in the Acronyms section at the end.
\usepackage[acronym,nomain,nopostdot]{glossaries}
\makenoidxglossaries
% Shared acronym macro container for EE800/820 lesson plan modules.
% Usage: type \IDE, \FPGA, \UART, etc. First use is expanded; later uses show acronym only.
% Requires in each document preamble:
%   \usepackage[acronym,toc]{glossaries}
%   \makeglossaries
%   % Shared acronym macro container for EE800/820 lesson plan modules.
% Usage: type \IDE, \FPGA, \UART, etc. First use is expanded; later uses show acronym only.
% Requires in each document preamble:
%   \usepackage[acronym,toc]{glossaries}
%   \makeglossaries
%   \input{../glossary}
% Print used acronyms near end of document:
%   \printglossary[type=\acronymtype,title={Acronyms Used}]

% Acronym definitions (only used entries appear in the printed glossary).
\newacronym{ask}{ASK}{Amplitude Shift Keying}
\newacronym{bga}{BGA}{Ball Grid Array}
\newacronym{bom}{BOM}{Bill of Materials}
\newacronym{bram}{BRAM}{Block Random Access Memory}
\newacronym{crc}{CRC}{Cyclic Redundancy Check}
\newacronym{css}{CSS}{Chirp Spread Spectrum}
\newacronym{dma}{DMA}{Direct Memory Access}
\newacronym{dsp}{DSP}{Digital Signal Processing}
\newacronym{eda}{EDA}{Exploratory Data Analysis}
\newacronym{fifo}{FIFO}{First In, First Out}
\newacronym{fpga}{FPGA}{Field-Programmable Gate Array}
\newacronym{fsk}{FSK}{Frequency Shift Keying}
\newacronym{fsm}{FSM}{Finite State Machine}
\newacronym{gpio}{GPIO}{General-Purpose Input/Output}
\newacronym{hal}{HAL}{Hardware Abstraction Layer}
\newacronym{i2c}{I\textsuperscript{2}C}{Inter-Integrated Circuit}
\newacronym{ide}{IDE}{Integrated Development Environment}
\newacronym{ila}{ILA}{Integrated Logic Analyzer}
\newacronym{ism}{ISM}{Industrial, Scientific, and Medical}
\newacronym{isr}{ISR}{Interrupt Service Routine}
\newacronym{lut}{LUT}{Look-Up Table}
\newacronym{mcu}{MCU}{Microcontroller Unit}
\newacronym{ml}{ML}{Machine Learning}
\newacronym{ook}{OOK}{On-Off Keying}
\newacronym{per}{PER}{Packet Error Rate}
\newacronym{rf}{RF}{Radio Frequency}
\newacronym{rssi}{RSSI}{Received Signal Strength Indicator}
\newacronym{rtl}{RTL}{Register-Transfer Level}
\newacronym{snr}{SNR}{Signal-to-Noise Ratio}
\newacronym{spi}{SPI}{Serial Peripheral Interface}
\newacronym{uart}{UART}{Universal Asynchronous Receiver/Transmitter}
\newacronym{vhdl}{VHDL}{VHSIC Hardware Description Language}

% Convenience wrappers so you can type \IDE, \RF, etc.
\providecommand{\ASK}{\gls{ask}}
\providecommand{\BGA}{\gls{bga}}
\providecommand{\BOM}{\gls{bom}}
\providecommand{\BRAM}{\gls{bram}}
\providecommand{\CRC}{\gls{crc}}
\providecommand{\CSS}{\gls{css}}
\providecommand{\DMA}{\gls{dma}}
\providecommand{\DSP}{\gls{dsp}}
\providecommand{\EDA}{\gls{eda}}
\providecommand{\FIFO}{\gls{fifo}}
\providecommand{\FPGA}{\gls{fpga}}
\providecommand{\FSK}{\gls{fsk}}
\providecommand{\FSM}{\gls{fsm}}
\providecommand{\GPIO}{\gls{gpio}}
\providecommand{\HAL}{\gls{hal}}
\providecommand{\IIC}{\gls{i2c}}
\providecommand{\IDE}{\gls{ide}}
\providecommand{\ILA}{\gls{ila}}
\providecommand{\ISM}{\gls{ism}}
\providecommand{\ISR}{\gls{isr}}
\providecommand{\LUT}{\gls{lut}}
\providecommand{\MCU}{\gls{mcu}}
\providecommand{\ML}{\gls{ml}}
\providecommand{\OOK}{\gls{ook}}
\providecommand{\PER}{\gls{per}}
\providecommand{\RF}{\gls{rf}}
\providecommand{\RSSI}{\gls{rssi}}
\providecommand{\RTL}{\gls{rtl}}
\providecommand{\SNR}{\gls{snr}}
\providecommand{\SPI}{\gls{spi}}
\providecommand{\UART}{\gls{uart}}
\providecommand{\VHDL}{\gls{vhdl}}


% Print used acronyms near end of document:
%   \printglossary[type=\acronymtype,title={Acronyms Used}]

% Acronym definitions (only used entries appear in the printed glossary).
\newacronym{ask}{ASK}{Amplitude Shift Keying}
\newacronym{bga}{BGA}{Ball Grid Array}
\newacronym{bom}{BOM}{Bill of Materials}
\newacronym{bram}{BRAM}{Block Random Access Memory}
\newacronym{crc}{CRC}{Cyclic Redundancy Check}
\newacronym{css}{CSS}{Chirp Spread Spectrum}
\newacronym{dma}{DMA}{Direct Memory Access}
\newacronym{dsp}{DSP}{Digital Signal Processing}
\newacronym{eda}{EDA}{Exploratory Data Analysis}
\newacronym{fifo}{FIFO}{First In, First Out}
\newacronym{fpga}{FPGA}{Field-Programmable Gate Array}
\newacronym{fsk}{FSK}{Frequency Shift Keying}
\newacronym{fsm}{FSM}{Finite State Machine}
\newacronym{gpio}{GPIO}{General-Purpose Input/Output}
\newacronym{hal}{HAL}{Hardware Abstraction Layer}
\newacronym{i2c}{I\textsuperscript{2}C}{Inter-Integrated Circuit}
\newacronym{ide}{IDE}{Integrated Development Environment}
\newacronym{ila}{ILA}{Integrated Logic Analyzer}
\newacronym{ism}{ISM}{Industrial, Scientific, and Medical}
\newacronym{isr}{ISR}{Interrupt Service Routine}
\newacronym{lut}{LUT}{Look-Up Table}
\newacronym{mcu}{MCU}{Microcontroller Unit}
\newacronym{ml}{ML}{Machine Learning}
\newacronym{ook}{OOK}{On-Off Keying}
\newacronym{per}{PER}{Packet Error Rate}
\newacronym{rf}{RF}{Radio Frequency}
\newacronym{rssi}{RSSI}{Received Signal Strength Indicator}
\newacronym{rtl}{RTL}{Register-Transfer Level}
\newacronym{snr}{SNR}{Signal-to-Noise Ratio}
\newacronym{spi}{SPI}{Serial Peripheral Interface}
\newacronym{uart}{UART}{Universal Asynchronous Receiver/Transmitter}
\newacronym{vhdl}{VHDL}{VHSIC Hardware Description Language}

% Convenience wrappers so you can type \IDE, \RF, etc.
\providecommand{\ASK}{\gls{ask}}
\providecommand{\BGA}{\gls{bga}}
\providecommand{\BOM}{\gls{bom}}
\providecommand{\BRAM}{\gls{bram}}
\providecommand{\CRC}{\gls{crc}}
\providecommand{\CSS}{\gls{css}}
\providecommand{\DMA}{\gls{dma}}
\providecommand{\DSP}{\gls{dsp}}
\providecommand{\EDA}{\gls{eda}}
\providecommand{\FIFO}{\gls{fifo}}
\providecommand{\FPGA}{\gls{fpga}}
\providecommand{\FSK}{\gls{fsk}}
\providecommand{\FSM}{\gls{fsm}}
\providecommand{\GPIO}{\gls{gpio}}
\providecommand{\HAL}{\gls{hal}}
\providecommand{\IIC}{\gls{i2c}}
\providecommand{\IDE}{\gls{ide}}
\providecommand{\ILA}{\gls{ila}}
\providecommand{\ISM}{\gls{ism}}
\providecommand{\ISR}{\gls{isr}}
\providecommand{\LUT}{\gls{lut}}
\providecommand{\MCU}{\gls{mcu}}
\providecommand{\ML}{\gls{ml}}
\providecommand{\OOK}{\gls{ook}}
\providecommand{\PER}{\gls{per}}
\providecommand{\RF}{\gls{rf}}
\providecommand{\RSSI}{\gls{rssi}}
\providecommand{\RTL}{\gls{rtl}}
\providecommand{\SNR}{\gls{snr}}
\providecommand{\SPI}{\gls{spi}}
\providecommand{\UART}{\gls{uart}}
\providecommand{\VHDL}{\gls{vhdl}}


% Call \printhandoutglossary just before \end{document} in each handout.
\newcommand{\printhandoutglossary}{%
	\clearpage
	\printnoidxglossary[type=\acronymtype,title={Glossary}]%
}

% ── Colors (match Module 1 lesson plan) ──────────────────────
\definecolor{stevens}{RGB}{163,31,52}
\definecolor{darkgray}{RGB}{60,60,60}
\definecolor{lightgray}{RGB}{240,240,240}
\definecolor{accentblue}{RGB}{30,80,140}

\hypersetup{
	colorlinks=true,
	linkcolor=stevens,
	urlcolor=accentblue,
	citecolor=darkgray,
}

% ── Defaults for header variables (handout overrides these) ──
\providecommand{\labid}{?}
\providecommand{\labtitleshort}{Lab Handout}

% ── Running header / footer ──────────────────────────────────
\pagestyle{fancy}
\fancyhf{}
\fancyhead[L]{\small\textcolor{darkgray}{Module 1 --- Lab Handout \labid}}
\fancyhead[R]{\small\textcolor{darkgray}{\labtitleshort}}
\fancyfoot[C]{\small\thepage}
\renewcommand{\headrulewidth}{0.4pt}

% ── Section formatting ───────────────────────────────────────
\titleformat{\section}{\Large\bfseries\color{stevens}}{\thesection}{0.6em}{}
\titleformat{\subsection}{\large\bfseries\color{darkgray}}{\thesubsection}{0.6em}{}

% ── Box styles ───────────────────────────────────────────────
\tcbset{
	infobox/.style={
		enhanced, colback=blue!3, colframe=accentblue, fonttitle=\bfseries,
		title={#1}, coltitle=white, breakable, left=6pt, right=6pt, top=4pt, bottom=4pt
	},
	warning/.style={
		enhanced, colback=red!3, colframe=stevens, fonttitle=\bfseries,
		title={#1}, coltitle=white, breakable, left=6pt, right=6pt, top=4pt, bottom=4pt
	},
	deliverable/.style={
		enhanced, colback=gray!5, colframe=darkgray, fonttitle=\bfseries,
		title={#1}, coltitle=white, breakable, left=6pt, right=6pt, top=4pt, bottom=4pt
	},
}

% ── Figure helper: always draws a placeholder box (images suppressed) ─
% Usage: \labfigure{filename}{width}{caption}{label}
\newcommand{\labfigure}[4]{%
	\begin{figure}[h!]
		\centering
		\fbox{\begin{minipage}[c][2.2in][c]{#2}
			\centering\color{darkgray}\small
				\textbf{[Figure placeholder]}\\[4pt]
				\texttt{\detokenize{#1}}\\[4pt]
				\textit{#3}
		\end{minipage}}%
		\caption{#3}
		\label{#4}
	\end{figure}%
}

% ── Title banner (call at top of document body) ──────────────
\newcommand{\labheader}[2]{%
	\begin{center}
		{\LARGE\bfseries\color{stevens} Lab Handout #1}\\[3pt]
		{\Large\bfseries #2}\\[6pt]
		{\small\color{darkgray} Capstone --- Module 1: \RF{} Hardware Foundations}
	\end{center}
	\vspace{4pt}
	{\color{darkgray}\rule{\textwidth}{0.4pt}}
	\vspace{10pt}
}


\begin{document}
\labheader{1-C}{\SPI{} Bring-Up: STM32 $\leftrightarrow$ SX1276 Register Access}

\section{Objective}
Establish reliable \SPI{} communication between the STM32 Nucleo-L476RG and the SX1276-based RFM95W module. You will read the \texttt{RegVersion} register and receive the expected value \texttt{0x12}, then demonstrate write/read-back and radio mode transitions. This handout is the first point at which firmware actually touches the \RF{} silicon.

\section{Prerequisites}
\begin{itemize}[leftmargin=2em]
	\item Lab Handout~1-A complete (breakout physically assembled).
	\item Lab Handout~1-B complete (toolchain working, blinky verified).
\end{itemize}

\section{Materials}
\begin{itemize}[leftmargin=2em]
	\item Assembled RFM95W breakout (from LH1-A).
	\item Nucleo-L476RG + USB cable (from LH1-B).
	\item Half-size breadboard, 8--10 male--female jumper wires.
	\item Logic analyzer (optional but strongly recommended for troubleshooting).
\end{itemize}

\section{Wiring}

\begin{table}[h!]
	\centering
	\begin{tabular}{l l l l}
		\toprule
		\textbf{RFM95W Pin} & \textbf{Function} & \textbf{Nucleo Pin} & \textbf{Arduino Hdr} \\
		\midrule
		\texttt{VIN} & +3.3\,V supply & \texttt{3V3} & 3V3 \\
		\texttt{GND} & Ground & \texttt{GND} & GND \\
		\texttt{SCK} & \SPI{} clock & \texttt{PA5} & D13 \\
		\texttt{MISO} & \SPI{} data in (to \MCU{}) & \texttt{PA6} & D12 \\
		\texttt{MOSI} & \SPI{} data out (from \MCU{}) & \texttt{PA7} & D11 \\
		\texttt{CS} & Chip select (active low) & \texttt{PB6} & D10 \\
		\texttt{RST} & Reset (active low) & \texttt{PA9} & D8 \\
		\texttt{G0} (\texttt{DIO0}) & Packet-done interrupt & \texttt{PA8} & D7 \\
		\bottomrule
	\end{tabular}
\end{table}

\begin{tcolorbox}[warning={Wire Before You Power}]
	Double-check \texttt{VIN} vs.\ \texttt{GND}. The RFM95W will tolerate 3.3--5\,V on \texttt{VIN} (it has an on-board regulator), but swapping \texttt{VIN} and \texttt{GND} \textit{will} destroy the module. Verify with a \mbox{DMM} before plugging the Nucleo into USB.
\end{tcolorbox}

\section{Procedure}

\subsection{Step 1 --- Configure peripherals via STM32CubeMX}
\begin{enumerate}[leftmargin=2em]
	\item Open the project from Lab Handout~1-B (or create a new one if you prefer to keep blinky separate). Open the CubeMX perspective.
	\item Enable \textbf{\SPI{}1} in \textbf{Full-Duplex Master} mode.
	\item Configure \SPI{}1 with the settings in Table~\ref{tab:spi-bringup}.
	\item Configure \texttt{PB6}, \texttt{PA9}, \texttt{PA8} as \mbox{GPIO} outputs named \texttt{RFM\_CS}, \texttt{RFM\_RST}, \texttt{RFM\_DIO0} (\texttt{DIO0} will become an \texttt{EXTI} input later; leave as plain \mbox{GPIO} input for this handout).
	\item Enable \textbf{USART2} in asynchronous mode at 115200\,8N1 (for \texttt{printf} redirection to the ST-LINK virtual COM port).
	\item Generate code.
\end{enumerate}

\begin{table}[h!]
	\centering
	\caption{\SPI{}1 configuration for RFM95W register access.}
	\label{tab:spi-bringup}
	\begin{tabular}{l l}
		\toprule
		\textbf{Parameter} & \textbf{Value} \\
		\midrule
		Mode & Full-Duplex Master \\
		Data size & 8 bits \\
		First bit & \mbox{MSB} first \\
		Prescaler & Set so that \SPI{} clock $\leq$5\,MHz (safe margin below the 10\,MHz max) \\
		\mbox{CPOL} & 0 (Low) \\
		\mbox{CPHA} & 1 Edge \\
		NSS & Software \\
		\bottomrule
	\end{tabular}
\end{table}

\subsection{Step 2 --- Redirect \texttt{printf} to USART2}
\begin{enumerate}[leftmargin=2em]
	\item In \texttt{main.c}, add an override for \texttt{\_\_io\_putchar} that calls \texttt{HAL\_UART\_Transmit(\&huart2, ...)}. A standard snippet:
\begin{verbatim}
int __io_putchar(int ch) {
    HAL_UART_Transmit(&huart2, (uint8_t*)&ch, 1, HAL_MAX_DELAY);
    return ch;
}
\end{verbatim}
	\item In the build configuration, ensure the \texttt{--specs=nano.specs} and \texttt{-u \_printf\_float} flags (or equivalent) are set so the reduced-newlib \texttt{printf} links.
\end{enumerate}

\subsection{Step 3 --- Write a minimal register-read function}
Implement the following helper:
\begin{verbatim}
uint8_t rfm_read_reg(uint8_t addr) {
    uint8_t tx[2] = { addr & 0x7F, 0x00 };  // bit 7 = 0 => read
    uint8_t rx[2] = { 0, 0 };
    HAL_GPIO_WritePin(RFM_CS_GPIO_Port, RFM_CS_Pin, GPIO_PIN_RESET);
    HAL_SPI_TransmitReceive(&hspi1, tx, rx, 2, HAL_MAX_DELAY);
    HAL_GPIO_WritePin(RFM_CS_GPIO_Port, RFM_CS_Pin, GPIO_PIN_SET);
    return rx[1];
}
\end{verbatim}
And the matching write helper (bit 7 = 1, followed by the value byte).

\subsection{Step 4 --- Reset the module and read \texttt{RegVersion}}
\begin{enumerate}[leftmargin=2em]
	\item In \texttt{main} after \HAL{} init, drive \texttt{RFM\_RST} low for at least 100\,$\mu$s, then high, then wait 10\,ms for the module to come out of reset.
	\item Call \texttt{rfm\_read\_reg(0x42)} and \texttt{printf} the result.
	\item Flash and run. Open the serial monitor (115200,8N1).
	\item You should see \texttt{RegVersion = 0x12}. Any other value indicates a wiring or configuration error.
\end{enumerate}

\begin{tcolorbox}[infobox={Why 0x12?}]
	The SX1276 silicon revision register returns \texttt{0x12} for the V1b mask set used in the RFM9x breakouts. This is the single most reliable sanity check for \SPI{} wiring, clock polarity, and reset sequencing: if you read \texttt{0x00} or \texttt{0xFF}, the module is not talking to you; if you read anything else, you have a clock phase or byte-order issue.
\end{tcolorbox}

\subsection{Step 5 --- Register write / read-back}
\begin{enumerate}[leftmargin=2em]
	\item Write the sequence \texttt{0xD9 0x0C 0xCC} to \texttt{RegFrfMsb}, \texttt{RegFrfMid}, \texttt{RegFrfLsb} (address \texttt{0x06}--\texttt{0x08}). This sets the carrier frequency to 906.5\,MHz. Note: the module must be in \texttt{SLEEP} mode before frequency writes take effect.
	\item Read the same registers back. Verify that all three values match what you wrote.
\end{enumerate}

\subsection{Step 6 --- Mode transitions}
\begin{enumerate}[leftmargin=2em]
	\item Read \texttt{RegOpMode} (\texttt{0x01}). Confirm the low 3 bits reflect the current mode.
	\item Write \texttt{0x80} to \texttt{RegOpMode} (LoRa mode, \texttt{SLEEP}).
	\item Write \texttt{0x81} (LoRa mode, \texttt{STANDBY}). Read back and confirm.
	\item Write \texttt{0x83} (LoRa mode, \texttt{TX}). Read back and confirm. The module will return to \texttt{STANDBY} after any transmission attempt; this is expected.
\end{enumerate}

\section{Verification Checklist}
\begin{itemize}[leftmargin=2em]
	\item[$\square$] \texttt{RegVersion} reads \texttt{0x12}.
	\item[$\square$] Write/read-back to \texttt{RegFrf*} registers returns identical values.
	\item[$\square$] \texttt{RegOpMode} mode bits transition correctly through \texttt{SLEEP}, \texttt{STANDBY}, \texttt{TX}.
	\item[$\square$] Serial output shows all register values in hex over the ST-LINK virtual COM port.
\end{itemize}

\begin{tcolorbox}[warning={Troubleshooting}]
	\begin{itemize}[leftmargin=1.2em]
		\item \textbf{\texttt{RegVersion} reads \texttt{0x00}}: \texttt{MISO} not connected, or \texttt{CS} never asserted.
		\item \textbf{\texttt{RegVersion} reads \texttt{0xFF}}: \texttt{MISO} floating (pull-up). Check the jumper to \texttt{PA6}.
		\item \textbf{Reads a seemingly-random but repeatable value}: \mbox{CPOL}/\mbox{CPHA} wrong, or bit 7 (read/write) flipped in your address byte.
		\item \textbf{First read is wrong, subsequent reads correct}: missing reset delay; increase the post-reset wait to 20\,ms.
	\end{itemize}
\end{tcolorbox}

\begin{tcolorbox}[deliverable={Lab Handout 1-C Deliverable}]
	\begin{enumerate}[leftmargin=1.2em]
		\item Source code for your \texttt{rfm\_read\_reg} / \texttt{rfm\_write\_reg} helpers and the \texttt{main} test sequence.
		\item Serial monitor log showing \texttt{RegVersion = 0x12}, frequency register write/read-back, and all three mode transitions.
		\item If a logic analyzer was used: one screenshot of a successful \texttt{RegVersion} read transaction with \texttt{CS}, \texttt{SCK}, \texttt{MOSI}, \texttt{MISO} traces.
	\end{enumerate}
\end{tcolorbox}

\end{document}
